\subsection{Real breeding case study: sorghum crop-model parameters}\label{sec:res-sorghum}
As an external check on private data of a different nature, we applied the
identical protocol to a sorghum panel of $136$ genotypes ($\sim$31{,}700 SNPs,
capped to $20{,}000$ by variance) in which the prediction targets are eight
parameters of an ecophysiological crop-growth model (e.g.\ radiation-use
efficiency and leaf-appearance rates) rather than measured phenotypes. Averaged
over the eight parameters under $5$-fold $\times\,3$-repeat cross-validation, the
affine-calibrated MLP trained with the Pearson loss raised predictive ability over
its MSE baseline by $\Delta r=+0.041$ (percentile-bootstrap CI $[0.004,\,0.076]$; $n=8$), with the hybrid
($+0.033$) and CCC ($+0.029$) losses close behind, and NDCG@$10$ improved by
$+0.018$ to $+0.021$. The effect size matches the public benchmark almost exactly,
although with only eight targets it is not individually significant (Holm
$p=0.33$). That an independent, proprietary dataset --- a different species, a
different target type, and a different research group --- reproduces the same
qualitative pattern is reassuring: aligning the neural training objective with the
correlation metric improves that metric, and affine calibration restores the
parameter scale. We report this as illustrative external evidence, not as an
additional powered test.
